% TOP_PROBLEM: {"id": 3093, "title": "Convex Equipartitions with Extreme Perimeter", "category_id": 6, "category": {"id": 6, "name": "geometry", "display_name": "Geometry", "description": "Euclidean and non-Euclidean geometry, geometric structures.", "slug": "geometry", "order_index": 6, "created_at": "2026-07-31T15:26:25.671Z"}, "status": "open", "problem_number": "OPG-60010", "set_id": 12, "statusQualification": "The region is specified as a compact convex body with interior, common boundaries are counted in each adjacent perimeter, and extrema are defined through infimum and supremum before attainment."}
% FIRST_POSTED: 2026-10-01
% ENTRY_KIND: statement_only
% UnsolvedMath ID 3093; source catalogue code OPG-60010
% !TeX program = lualatex
% Definition and sources only; this file is not an LLM solution attempt.
\documentclass[11pt,a4paper]{article}
\usepackage[margin=25mm]{geometry}
\usepackage{fontspec}
\setmainfont{lmroman10-regular.otf}[BoldFont=lmroman10-bold.otf,
 ItalicFont=lmroman10-italic.otf,BoldItalicFont=lmroman10-bolditalic.otf]
\usepackage{amsmath,amssymb,mathtools}
\usepackage{unicode-math}
\setmathfont{latinmodern-math.otf}
\AtBeginDocument{\let\setminus\smallsetminus}
\usepackage{xurl,hyperref}
\hypersetup{colorlinks=true,urlcolor=blue}
\setcounter{secnumdepth}{0}
\setlength{\parindent}{0pt}
\setlength{\parskip}{5pt}
\setlength{\emergencystretch}{3em}
\begin{document}
\section*{Convex Equipartitions with Extreme Perimeter}\label{3093}
{\small UnsolvedMath ID 3093 · OPG-60010 · Geometry · OpenGarden}
\subsection{Definitions and mathematical statement}
Let $C\subset\mathbb R^2$ be a compact convex set with nonempty interior,
and let $n\ge1$ be an integer. Write $|C|$ for its area and
$\operatorname{per}(C)$ for the Euclidean length of its boundary.
An admissible convex equipartition is a family of compact convex sets
$C_1,\ldots,C_n$ whose union is $C$, whose interiors are pairwise disjoint,
and for which $|C_i|=|C|/n$ for every $i$.

The quantity to optimize is the sum
\[
 P(C_1,\ldots,C_n)=\sum_{i=1}^n\operatorname{per}(C_i).
\]
Determine its least and greatest possible values over admissible partitions,
and describe partitions attaining them. Equivalently one can first define
$P_{\min}(C,n)$ and $P_{\max}(C,n)$ as the infimum and supremum and
include attainment in a complete determination. For a polygon the pieces
are understood as closed regions meeting along their cut boundaries,
with boundary overlaps allowed but no overlap of positive area.

The exterior boundary of $C$ is counted once in the sum, while every
internal interface of positive length is counted twice. Thus, writing
$L$ for the total length of internal interfaces counted once,
$P=\operatorname{per}(C)+2L$. The optimization is also optimization of
cut length. Only areas must be equal; piece perimeters may differ and
forcing them to coincide would change the feasible family.

The question ranges over the input body and number of pieces, including
nonsymmetric bodies. Merely establishing that equal-area convex partitions
exist does not identify either perimeter extremum.
\subsection{Short English statement}
Among convex partitions into equally large areas, which arrangements minimize or maximize the sum of piece perimeters?
\subsection{Sources}
\begin{itemize}
\item[S1] ULAM.AI and the UnsolvedMath Contributors, supplied problems.json, record 3093 (OPG-60010), OpenGarden collection. Catalogue identity and source transcription; CC BY 4.0.
\url{https://huggingface.co/datasets/ulamai/UnsolvedMath}.
\item[S2] Open Problem Garden, Convex Equipartitions with Extreme Perimeter. Original problem and discussion; the mathematical formulation below is expanded from the supplied source record.
\url{http://www.openproblemgarden.org/op/convex_equipartitions_with_extreme_perimeter}.
\end{itemize}
\end{document}
